\documentclass[11pt,a4paper]{article}

\usepackage[margin=1in]{geometry}
\usepackage{hyperref}
\usepackage{enumitem}
\usepackage{booktabs}
\usepackage{xcolor}

\definecolor{tier1}{RGB}{180,40,40}
\definecolor{tier2}{RGB}{40,100,160}
\definecolor{tier3}{RGB}{40,140,80}
\definecolor{tier4}{RGB}{120,120,120}

\title{Reading List for the Algebraic Islands Project\\[4pt]
\large Prioritized by relevance to current research}
\author{Michael Kernaghan\\Pacific Quantum Systems}
\date{February 2026}

\begin{document}

\maketitle

\noindent This reading list supports the paper ``The Algebraic Landscape of Kochen--Specker Sets in Dimension Three'' and its companion PRL letters.  Papers are ordered by priority within four tiers.  Current results to keep in mind while reading: OCUS proof that 31 is optimal in the 49-ray integer pool; universal merge saturation across all six islands (3{,}756 merges); $N(S) = 6$ for the Eisenstein island (confirming the Cortez--Morales--Reyes minimal ring); SI-C closure convergence with the integer pool.

\bigskip

%% ============================================================
\section*{\textcolor{tier1}{Tier 1 --- Read Carefully}}
\addcontentsline{toc}{section}{Tier 1 --- Read Carefully}

These papers directly shape the arguments and claims in our work.

\begin{enumerate}[label=\textbf{\arabic*.}, leftmargin=2em]

\item \textbf{S.~Trandafir and A.~Cabello},
``Two fundamental solutions to the rigid Kochen--Specker set problem and the solution to the minimal Kochen--Specker set problem under one assumption,''
\textit{Phys.\ Rev.\ A} \textbf{111}, 052204 (2025).
\href{https://arxiv.org/abs/2501.11640}{arXiv:2501.11640}.

\smallskip\noindent
\textit{Why}: The single most relevant paper.  Proves CK-31 is rigid (unique up to unitary equivalence).  Constructs the 97-element SI-C closure---our computation shows it overlaps our integer pool in 37 rays and contains it after a second closure round.  Introduces bootstrap percolation for KS sets (Proposition~1).  Their Conjecture~1 (no rigid KS set $\leq 30$ containing the SI-C set) is complementary to our OCUS proof.  Re-read with the SI-C closure and $N(S)$ results fresh.

\bigskip

\item \textbf{V.~Cortez, R.~Morales, and E.~J.~Reyes},
``Minimal ring extensions of the integers exhibiting Kochen--Specker contextuality,''
\href{https://arxiv.org/abs/2211.13216}{arXiv:2211.13216} (2022).

\smallskip\noindent
\textit{Why}: Proves $\mathbb{Z}[1/6]$ is the minimal ring extension of $\mathbb{Z}$ exhibiting KS contextuality in dimension~3.  Their $N(S) = \operatorname{lcm}\{\|v\|^2 : v \in S\}$ invariant is the bridge between their ring-theoretic framework and our algebraic islands: our computation confirms the Eisenstein island achieves exactly $N = 6$, the minimal value.  Understanding their proof technique (Theorem~1: $6 \mid N$ is necessary and sufficient) could help prove the $6 \mid n$ conjecture for cyclotomic fields.

\bigskip

\item \textbf{Z.~Li, C.~Bright, and V.~Ganesh},
``Leveraging SAT solvers for the minimum Kochen--Specker problem,''
\textit{Proceedings of IJCAI-24}, pp.~1944--1952 (2024).
\href{https://arxiv.org/abs/2306.13319}{arXiv:2306.13319}.

\smallskip\noindent
\textit{Why}: Establishes the current lower bound of 24 vectors via SAT+CAS orderly generation.  Their pipeline (enumerate abstract KS hypergraphs $\to$ check realizability via Z3) is directly complementary to our finite-pool SAT encoding (C4a--C4c), which could serve as a ``quick accept'' screen.  Key details: 17 known minimal unembeddable subgraphs, 41 KS candidates at $n = 23$ (all unembeddable), estimated $\sim$125 CPU-years to push to $n = 24$.  Read carefully if pursuing the LBG filter avenue.

\bigskip

\item \textbf{M.~Kirchweger, A.~Peitl, and S.~Szeider},
``Co-certificate learning with SAT modulo symmetries,''
\textit{Proceedings of AAAI-23}, pp.~4119--4126 (2023).

\smallskip\noindent
\textit{Why}: Independent lower bound of 24.  Co-certificate learning (CCL) is conceptually related to our OCUS/CEGAR blocking---both learn structural constraints from satisfying assignments.  Key data: 3{,}747{,}950 KS candidates at $n = 23$, of which only 2 survive their unembeddable-subgraph filter.  The explosive growth ($88{,}282$ at $n = 22 \to 3.7$M at $n = 23$) shows why pure enumeration cannot reach $n = 31$.

\end{enumerate}

\bigskip

%% ============================================================
\section*{\textcolor{tier2}{Tier 2 --- Read for Depth}}
\addcontentsline{toc}{section}{Tier 2 --- Read for Depth}

These strengthen connections already made in our papers.

\begin{enumerate}[label=\textbf{\arabic*.}, leftmargin=2em, start=5]

\item \textbf{A.~Cabello},
``Simplest bipartite perfect quantum strategies,''
\textit{Phys.\ Rev.\ Lett.} \textbf{134}, 010201 (2025).
\href{https://arxiv.org/abs/2311.17735}{arXiv:2311.17735}.

\smallskip\noindent
\textit{Why}: Definition~2 (B-KS) is the foundation of our BPQS analysis.  Theorem~1 (every BPQS defines a KS set) gives operational significance to all six islands.  A careful re-read ensures our encoding of conditions (I$'$) and (II$'$) is airtight.

\bigskip

\item \textbf{A.~Cabello},
``Simplest Kochen--Specker set,''
\textit{Phys.\ Rev.\ Lett.} (2025).
\href{https://arxiv.org/abs/2508.07335}{arXiv:2508.07335}.

\smallskip\noindent
\textit{Why}: The 33-vector, 14-basis Eisenstein construction via Weyl--Heisenberg.  Table~I has BPQS numbers we independently verify.  This is the concrete bridge between our Eisenstein island and Cabello's ``simplest'' criterion.  His coordinates $\{0, 1, \omega, \bar\omega\}$ are precisely our Eisenstein alphabet.

\bigskip

\item \textbf{Y.~Liu, H.~Y.~Chung, E.~Z.~Cruzeiro, J.~R.~Gonzales-Ureta, R.~Ramanathan, and A.~Cabello},
``Equivalence between face nonsignaling correlations, full nonlocality, all-versus-nothing proofs, and pseudotelepathy,''
\textit{Phys.\ Rev.\ Research} \textbf{6}, L042035 (2024).
\href{https://arxiv.org/abs/2310.10600}{arXiv:2310.10600}.

\smallskip\noindent
\textit{Why}: The chain FNS $=$ FN $=$ AVN $=$ PT, all equivalent to KS sets.  This is what makes our island classification operationally significant: each island defines not just a KS set but an entire class of perfect quantum strategies.

\end{enumerate}

\bigskip

%% ============================================================
\section*{\textcolor{tier3}{Tier 3 --- Read for New Avenues}}
\addcontentsline{toc}{section}{Tier 3 --- Read for New Avenues}

These open new directions of research.

\begin{enumerate}[label=\textbf{\arabic*.}, leftmargin=2em, start=8]

\item \textbf{S.~Abramsky, R.~S.~Barbosa, and S.~Mansfield},
``The contextual fraction as a measure of contextuality,''
\textit{Phys.\ Rev.\ Lett.} \textbf{119}, 050504 (2017).
\href{https://arxiv.org/abs/1705.07918}{arXiv:1705.07918}.

\smallskip\noindent
\textit{Why}: The contextual fraction is a different measure from CSW---it quantifies the ``amount'' of contextuality via a linear program.  Computing it for each island would give a third axis (alongside CSW and BPQS) for the Pareto frontier across algebraic substrates.

\bigskip

\item \textbf{C.~Budroni, A.~Cabello, O.~G\"uhne, M.~Kleinmann, and J.-\AA.~Larsson},
``Kochen--Specker contextuality,''
\textit{Rev.\ Mod.\ Phys.} \textbf{94}, 045007 (2022).
\href{https://arxiv.org/abs/2102.13036}{arXiv:2102.13036}.

\smallskip\noindent
\textit{Why}: The comprehensive review ($\sim$100 pages).  Essential for ensuring we haven't missed a connection and for framing the introduction and related work.  Particularly relevant: Section~IV (dimension-specific results), Section~VI (operational significance), and the open problems list.

\bigskip

\item \textbf{A.~Peres},
\textit{Quantum Theory: Concepts and Methods} (Kluwer, 1993), Chapter~7.

\smallskip\noindent
\textit{Why}: The original source for CK-31, with cube diagrams showing how the coordinate alphabet $\{0, \pm 1, \pm 2\}$ generates the construction.  Worth revisiting to ensure our alphabet-based narrative is faithful to the original presentation.

\end{enumerate}

\bigskip

%% ============================================================
\section*{\textcolor{tier4}{Tier 4 --- Skim for Awareness}}
\addcontentsline{toc}{section}{Tier 4 --- Skim for Awareness}

Background and historical context.

\begin{enumerate}[label=\textbf{\arabic*.}, leftmargin=2em, start=11]

\item \textbf{M.~Pavicic, N.~D.~Megill, B.~C.~Waegell, and P.~K.~Aravind},
``Automated generation of Kochen--Specker sets,''
\textit{Sci.\ Rep.} \textbf{9}, 6765 (2019).

\smallskip\noindent
\textit{Why}: Their automated generation approach and ``vector component'' terminology (our ``coordinate alphabet'').  Relevant mainly to avoid the terminological dispute (Pavicic distinguishes ``KS sets'' from ``extended KS sets''; we use standard definitions).

\bigskip

\item \textbf{S.~Uijlen and B.~Westerbaan},
``A Kochen--Specker system has at least 22 vectors,''
\textit{New Generation Computing} \textbf{34}, 3--23 (2016).

\smallskip\noindent
\textit{Why}: Historical context only---the previous lower bound (22), now superseded by the 24 of Li--Bright--Ganesh and Kirchweger et al.

\end{enumerate}

\bigskip

%% ============================================================
\section*{Reading Strategy}

\begin{itemize}[leftmargin=1.5em]
\item Papers 1--4 deserve careful, line-by-line reading.  Budget 2--3 hours each.
\item Papers 5--7 are quick re-reads to verify specific claims.  Budget 30--60 minutes each.
\item Papers 8--10 are for whichever new avenue you choose to pursue.
\item Papers 11--12 are 15-minute skims.
\end{itemize}

\medskip
\noindent\textit{Total estimated reading time: 15--20 hours.}

\end{document}
